\section{結果}
\label{sec:results}
\subsection{比較手法との比較}
\label{sec:prior}
表~\ref{tab:main}に test での結果を示す. \Ours{}の5つの平均は\olTestOL\%で,先行研究の最良(\olBestPriorName, \olBestPriorMean\%)を\olGapCfsKnn{}ポイント上回る([\olGapCfsKnnLo, \olGapCfsKnnHi]). 差が大きいのはCIFAR-10とImageNet-64である(\olTestCten\%と\olTestIN\%. 先行研究の最良は\olPriorBestCten\%と\olPriorBestIN\%). 上下反転した分布内の画像も\olTestVf\%で見分け,先行研究の最良(\olVfBestPriorAvg\%)を上回る. 及ばないのはSVHNで(\olTestSVHN\%. \olPriorBestSVHNName{}は\olPriorBestSVHN\%),全\olPairs{}組のうち先行研究のどれかに負けた\olNLostPrior{}組は,すべてSVHNを分布内とする組である. 相手は10回以上の評価を要する尤度やスコアの大きさの手法で, CFSの2通りに負けた組はない.

\begin{table}[t]\centering\scriptsize\setlength{\tabcolsep}{2.2pt}\renewcommand{\arraystretch}{1.0}
\caption{時間と用意するもの(CIFAR-10のモデル, CPU 2スレッド). 止める：順伝播を止める位置(1/2：デコーダの1/2の解像度). 比：ブロックAまでの順伝播(\olFwdA{}秒)との比. 画像：事前に通す分布内の画像. 数値, MB：保存する数値の数と16ビットでの大きさ. AUROC：test での5つの平均.}
\label{tab:cost}
\begin{adjustbox}{max width=\columnwidth}
\begin{tabular}{@{}lccccccc@{}}\toprule
手法 & 止める & 秒/枚 & 比 & 画像 & 数値 & MB & AUROC\\\midrule
\inputrows{tab_cost_ol.tex}
\end{tabular}
\end{adjustbox}
\end{table}
\begin{table*}[t]\centering\scriptsize\setlength{\tabcolsep}{3pt}\renewcommand{\arraystretch}{0.92}
\caption{val2 での分析(AUROC, \%. ノイズの種3つの平均). 平均：自然画像の分布外集合に対する5つの平均. 対照：縮小の履歴をそろえた対照(\ref{sec:hist}節)での平均. 上下反転：上下反転した分布内の画像. $c=8$ と7つの向きは,この分割での開発で決めた(補足資料A.3). 分布内ごとの値は補足資料の表~\ref{tab:devfull}.}
\label{tab:dev}
\begin{tabular}{@{}lccc|lccc|lccc@{}}\toprule
\multicolumn{4}{@{}c|}{定数 $c$(7つの向き)} & \multicolumn{4}{c|}{使う向き($c=8$)} & \multicolumn{4}{c@{}}{分布内の画像の枚数 $n$}\\
 & 平均 & 対照 & 上下反転 & & 平均 & 対照 & 上下反転 & & 平均 & 対照 & 上下反転\\\midrule
\inputrows{tab_devw_ol.tex}
\end{tabular}
\end{table*}
\begin{table*}[t]\centering\scriptsize\setlength{\tabcolsep}{3.4pt}\renewcommand{\arraystretch}{0.92}
\caption{$256{\times}256$ の結果(AUROC, \%. 分布内はImageNet,各集合1{,}000枚). 下段は112ピクセルを経由した対照. 近い：SSB-hardとNINCO. 遠い：ほかの3つ. \colorbox{firstblue}{\textbf{水色の太字}}：\Ours{}がCFSの2通りに並ぶか上回る列.}
\label{tab:hr}
\begin{adjustbox}{max width=\textwidth}
\begin{tabular}{@{}llccccc|cc|c|c@{}}\toprule
画像 & 手法 & SSB-hard & NINCO & iNaturalist & Textures & OpenImage-O & 近い & 遠い & 5つの平均 & 上下反転\\\midrule
\inputrows{tab_hr_ol.tex}
\end{tabular}
\end{adjustbox}
\end{table*}
\subsection{6つの分割}
\label{sec:six}
表~\ref{tab:six}のとおり, 5つの平均は\olSixMin{}〜\olSixMax\%(6分割の平均\olSixMean\%)である. 上下反転した分布内の画像は\olSixVfMin{}〜\olSixVfMax\%で,分布内ごとの差が大きい(CelebA \olSixVfCelebA\%, SVHN \olSixVfSVHN\%). 左右反転した自然画像は分布内の画像なので,見分けない(ほぼ50\%).

\subsection{時間と用意するもの}
\label{sec:cost}
表~\ref{tab:cost}のとおり, \Ours{}は1枚\olSecOL{}秒で,その大半はブロックAまでの順伝播である. 途中で止めたCFSの\olSecRatioCfs{}倍(対角ガウス)と\olSecRatioCfsKnn{}倍(kNN ヘッド)で,保存するのは kNN ヘッドの\olMBCfsKnn\,MBに対して\olMBOurs\,MBである. CFSを\Ours{}と同じ val1 の画像に当てはめても, test での値は\olCfsValMean\%(対角ガウス)と\olCfsKnnValMean\%(kNN ヘッド)で,差は画像の枚数によらない.

\subsection{何が効いているか}
\label{sec:ablate}
表~\ref{tab:six}の下の3行のとおり,距離の項だけでは\olSixT\%,内積の項だけでは\olSixL\%で,内積の項が主である. 距離の項はSVHNとImageNet-64で効く. 全ブロックを使っても平均は\olSixAll\%でほぼ同じだが, SVHNが上がり, CIFAR-10とImageNet-64が下がる. val2 での分析(表~\ref{tab:dev})では, $c$ を0から大きくすると\olDevCzero\%から上がり, 8と16で最も高く(\olDevCeight\%と\olDevCsixteen\%), $c\to\infty$ で\olDevCinf\%に下がる. 上下反転だけを使うと\olDevOv\%で, SVHNが低い. 分布内の画像を1{,}000枚に減らしても\olDevNthousand\%である. 定数のない形(8つの向きのうち「そのまま」である事後確率)は\olDevPk\%にとどまる.

\subsection{256{$\times$}256の画像}
\label{sec:hr}
\OLHRTEXT

\section{検証：何を見ているか}
\label{sec:verify}
\begin{table}[t]\centering\scriptsize\setlength{\tabcolsep}{2.2pt}\renewcommand{\arraystretch}{0.95}
\caption{振幅は下がるか(val2,ブロックAまで). 下がる,上がる： $z<-0.2$, $z>0.2$ のブロックの割合(\%). $z_{\mathrm{A}}$：ブロックAでの $z$. 分布外の集合は集合ごとの値の平均.}
\label{tab:fall}
\begin{adjustbox}{max width=\columnwidth}
\begin{tabular}{@{}lc|ccc|ccc|ccc@{}}\toprule
 & & \multicolumn{3}{c|}{向きを変えた画像} & \multicolumn{3}{c|}{分布外の集合} & \multicolumn{3}{c}{val2 の分布内}\\
分布内 & $p$ & 下がる & 上がる & $z_{\mathrm{A}}$ & 下がる & 上がる & $z_{\mathrm{A}}$ & 下がる & 上がる & $z_{\mathrm{A}}$\\\midrule
\inputrows{tab_fall_ol.tex}
\end{tabular}
\end{adjustbox}
\end{table}
\begin{table}[t]\centering\scriptsize\setlength{\tabcolsep}{2.4pt}\renewcommand{\arraystretch}{0.95}
\caption{分布内の画像と劣化させた写しを見分けるAUROC(\%. test6, 5つの平均. 50\%なら見分けない). ぼかしは標準偏差(ピクセル), JPEGは品質,縮小は倍率(縮小して戻す). 分布内ごとの値は補足資料の表~\ref{tab:degrid}.}
\label{tab:degr}
\begin{adjustbox}{max width=\columnwidth}
\begin{tabular}{@{}lccccccc|c@{}}\toprule
手法 & \multicolumn{2}{c}{ぼかし} & \multicolumn{3}{c}{JPEG} & \multicolumn{2}{c|}{縮小} & 平均\\
 & 0.5 & 1.0 & 90 & 75 & 50 & 3/4 & 1/2 & \\\midrule
\inputrows{tab_degr_ol.tex}
\end{tabular}
\end{adjustbox}
\end{table}
\begin{table*}[t]\centering\scriptsize\setlength{\tabcolsep}{3.4pt}\renewcommand{\arraystretch}{0.92}
\caption{縮小の履歴をそろえた対照(AUROC, \%. test). 対照では,すべての画像の一辺を半分に縮小して戻す. 右端：\Ours{}との差と95\%信頼区間.}
\label{tab:hist}
\begin{adjustbox}{max width=\textwidth}
\begin{tabular}{@{}llccccc|c|c|c@{}}\toprule
画像 & 手法 & CIFAR-10 & SVHN & CelebA & ImageNet-64 & FashionMNIST & 5つの平均 & 上下反転 & \Ours{}との差\\\midrule
\inputrows{tab_hist_ol.tex}
\end{tabular}
\end{adjustbox}
\end{table*}
\begin{table}[t]\centering\scriptsize\setlength{\tabcolsep}{2.2pt}\renewcommand{\arraystretch}{0.95}
\caption{重みを変えたときのAUROC(\%. test6,ノイズの種0). 較正の画像と式は同じ. 初期化したまま：学習前の重み. 並べ替え：テンソルごとに要素を並べ替えた重み. 上下：上下反転した分布内の画像. 内積,距離：内積の項だけ,距離の項だけでの5つの平均.}
\label{tab:var}
\begin{adjustbox}{max width=\columnwidth}
\begin{tabular}{@{}lccccc|c|c|cc@{}}\toprule
重み & C-10 & SVHN & CelebA & IN-64 & FMNIST & 平均 & 上下 & 内積 & 距離\\\midrule
\inputrows{tab_var_ol.tex}
\end{tabular}
\end{adjustbox}
\end{table}
\subsection{振幅は下がるか}
\label{sec:fall}
ブロックごとに,集合の振幅の平均と分布内の平均の差を分布内の標準偏差で割った値を $z$ とし, $z<-0.2$ を「下がる」と数えた(符号だけで数えると val2 の分布内の画像でも偏る. 補足資料の表~\ref{tab:fallsign}). 表~\ref{tab:fall}のとおり,向きを変えるとブロックAまでの\olFallOrMin{}〜\olFallOrMax\%のブロックで振幅が下がり,上がるのは\olRiseOrMax\%以下で,ブロックAでは5つすべてで下がる. 分布外の集合でも, \olNSetsA{}集合すべてで,ブロックAの振幅の平均が分布内より低い. ただしSVHNでは,分布外の集合で上がるブロックのほうが多い. val2 の分布内の画像では,下がるブロックも上がるブロックもない.

\subsection{画質への感度}
\label{sec:qual}
内容が同じでも, \Ours{}は劣化させた写しを平均\olDqOLMean\%で見分ける(表~\ref{tab:degr}). CFSの kNN ヘッドも\olDqCfsKnnMean\%で,画質への感度は中間特徴を読む検知器に共通である. \Ours{}では距離の項が内積の項より敏感で, ImageNet-64は品質90のJPEGを\olDqINJpeg\%で見分ける.

\subsection{縮小の履歴をそろえた対照}
\label{sec:hist}
分布外集合は,データセットごとに違う元画像から違う方法で縮小されている. すべての画像(較正とCFSの参照画像を含む)の一辺を半分に縮小して戻すと, \Ours{}は\olEqDropOL{}ポイント, CFSの kNN ヘッドも\olEqDropCfsKnn{}ポイント下がり,両者の差は\olEqGapCfsKnn{}ポイント([\olEqGapCfsKnnLo, \olEqGapCfsKnnHi])で変わらない(表~\ref{tab:hist}). 6つの分割では\olEqSixDrop{}ポイント下がり(補足資料の表~\ref{tab:splitseq}),主にImageNet-64(\olSixIN\%から\olEqSixIN\%)とCIFAR-10による. この低下には,履歴の違いによる分と細部が失われた分が混ざっている.

\subsection{モデルの学習}
\label{sec:model}
ノイズ除去器の重みを初期化したままにすると, 5つの平均は\olVarTrained\%から\olVarRinit\%に,要素を並べ替えると\olVarRshuf\%に下がる(test6. 表~\ref{tab:var}). 内積の項は\olVarTrainedL\%から\olVarRinitL\%に下がり,学習した応答を読んでいる. 距離の項の低下は小さく(\olVarTrainedT\%から\olVarRinitT\%),画像の統計の違いは学習していないネットワークの振幅にも現れる. 別のデータで学習したモデルでも,数字と自然画像は距離の項で分かれる(SVHNの画像をCIFAR-10のモデルに通すと\olVarByMax\%). SVHNの結果は,モデルが数字を学習した証拠にはならない.

\subsection{分布外集合から決めた方向との比較}
\label{sec:oracle}
val2 の分布外集合を使って決めた線形判別の方向(手法ではない)と比べると, \Ours{}の方向 $\Pm(\vmu-\vmu_{o})$ との余弦は\olCosMin{}〜\olCosMax{}で,分布外の画像なしで近い方向が得られている(補足資料の表~\ref{tab:oracle}). 5つの平均は\Ours{}の\olXOLR\%に対し,まとめて決めた方向で\olXPool\%,集合ごとに決めた方向で\olXSet\%である. ただし後者は画質や履歴の違いも拾う(ImageNet-64に対するSSB-hardを\olXSsbSet\%で分ける. \Ours{}は\olXSsbOLR\%).

\section{考察}
\label{sec:disc}
\textbf{何をしているか.} \Ours{}は,補助の分布外データ~\cite{hendrycks2019oe}の代わりに分布内の画像の向きを変えた写しを使い,学習の代わりに平均を1つ求める. 向きを変えた画像でも分布外の画像でもブロックAの振幅は下がり(\ref{sec:fall}節),得られる方向は分布外集合で決めた方向に近い(\ref{sec:oracle}節).

\textbf{限界.} (1) $c=8$ は val2 で決めた(8と16の間は平ら). $c$ は尤度比の共分散の比 $1+1/c$ にあたるが,測った比とは合わない. (2) 開発では候補の式を\olNDev{}通り試し,変種を\olNVar{}通り計算した. テスト分割はどれも以前に他の手法で結果を見ている. (3) 縮小の履歴と画質に反応する(CFSも同じ). (4) 向きを変えても並びが動かないデータでは内積の項が働かない. $256{\times}256$ のImageNetがそうで,内容を見分けている証拠も得られなかった. 上下反転した数字も見分けにくい(SVHN \olSixVfSVHN\%). (5) SVHNでは尤度やスコアの大きさを使う手法に及ばず,内容が近い集合(ImageNet-64)は難しい. 分類器の特徴を使う検知器とは比べていない.

\section{おわりに}
凍結した拡散モデルのブロックごとに数値を1つ読み,向きを変えた分布内の画像から求めた方向で判定する\Ours{}を提案した. 5つの分布内データセットで,平均AUROCは test で\olTestOL\%, 6つの分割で\olSixMin{}〜\olSixMax\%で,先行研究の最良を\olGapCfsKnn{}ポイント上回った. 必要なのは,途中で止める順伝播1回と\olNStored{}個の数値である. 一方,向きを変えても並びが動かない場合($256{\times}256$ のImageNet)には働かない. $\Delta$ は分布内の画像だけから求まるので,使えるかどうかは分布外の画像なしで確かめられる.
